\par\textbf{Data}\par

\par\textbf{Twitter Data (United Kingdom)}\par

\begin{itemize}
\item
  Geocoded tweets mentioning Israel, Jews or Tottenham (and variations thereof) from 2013 to 2022 are aggregated to 355 districts per day.\footnote{(Binstok, Gould, and Kaplan 2024, p. 2335, l. 16--22)}
\item
  Days with military hostilities between Israel and the Palestinians are excluded, leaving 1,253,505 district-days.\footnote{(Binstok, Gould, and Kaplan 2024, p. 2335, l. 22--25)}
\item
  Sentiment is classified with TweetNLP (a RoBERTa model), which assigns each tweet a probability of being positive, neutral or negative.\footnote{(Binstok, Gould, and Kaplan 2024, p. 2335, l. 26--29; p. 2336, l. 23--24)}
\item
  The negative count of a district-day is the sum of the tweets' probabilities of being negative (e.g. 0.4 + 0.4 + 0.6 = 1.4).\footnote{(Binstok, Gould, and Kaplan 2024, p. 2336, l. 24--31)}
\item
  Most tweets are neutral; tweets about Jews are twice as often negative as positive, tweets about Israel almost four times as often.\footnote{(Binstok, Gould, and Kaplan 2024, p. 2336, l. 32--38 (Table 1))}
\end{itemize}

\par\textbf{German General Social Survey (ALLBUS 2016)}\par

\begin{itemize}
\item
  Source: GESIS -- Leibniz-Institut für Sozialwissenschaften, study ZA5252, doi:10.4232/1.12837.\footnote{(Binstok, Gould, and Kaplan 2024, fn. 16, p. 2344, l. 36--38)}
\item
  ALLBUS is the only official survey that contains both standard antisemitism items and questions on Israel; the Israel items were included only in the 2016 wave.\footnote{(Binstok, Gould, and Kaplan 2024, p. 2333, l. 21--24; p. 2345, l. 2--7)}
\item
  \textbf{Sample:} respondents born in Germany between 1930 and 1990 who are Christian or have no religious affiliation, thereby excluding Jews and Muslims.\footnote{(Binstok, Gould, and Kaplan 2024, p. 2345, l. 8--12)}

  \begin{itemize}
  \item
    The final sample comprises 2,336 respondents with non-missing values for all main variables.\footnote{(Binstok, Gould, and Kaplan 2024, notes to Table 6, p. 2347, l. 34--38)}
  \end{itemize}
\item
  \textbf{Outcome (Y):} agreement on a scale from 1 to 7 with the statement ``It is unjust that Israel is taking away land from the Palestinians''.\footnote{(Binstok, Gould, and Kaplan 2024, p. 2345, l. 12--17)}

  \begin{itemize}
  \item
    Second outcome, used for robustness: ``Given Israel's policies, I can well understand that people are against the Jews''.\footnote{(Binstok, Gould, and Kaplan 2024, p. 2345, l. 14--16)}
  \item
    The first outcome serves as the main outcome because, unlike the second, it does not refer to Jews.\footnote{(Binstok, Gould, and Kaplan 2024, p. 2346, l. 13--20)}
  \item
    Its mean is 5.30 (standard deviation 1.62), i.e. the distribution is skewed towards agreement.\footnote{(Binstok, Gould, and Kaplan 2024, p. 2346, l. 2--3; Table 6, p. 2347, l. 10)}
  \item
    The item is ambiguous: disagreement can mean that the taking of land is considered just or that no land is being taken.\footnote{(Binstok, Gould, and Kaplan 2024, p. 2346, l. 4--6)}
  \end{itemize}
\item
  \textbf{Variable of interest (X):} agreement on a scale from 1 to 7 with the statement ``Jews have too much influence in the world''.\footnote{(Binstok, Gould, and Kaplan 2024, p. 2348, l. 7--8; p. 2349, l. 2--4)}

  \begin{itemize}
  \item
    All five antisemitism items are coded such that higher values indicate more negative views about Jews.\footnote{(Binstok, Gould, and Kaplan 2024, p. 2348, l. 2--13)}
  \item
    For robustness, the mean of the five items and their first factor are used (factor loadings 0.67, 0.62, 0.65, 0.28 and 0.43).\footnote{(Binstok, Gould, and Kaplan 2024, p. 2349, l. 4--8 and fn. 18, l. 19--21)}
  \item
    The mean of the main item is 3.09 (standard deviation 1.92).\footnote{(Binstok, Gould, and Kaplan 2024, Table 6, p. 2347, l. 12)}
  \end{itemize}
\item
  \textbf{Instrument (Z):} an indicator for being born between 1930 and 1945.\footnote{(Binstok, Gould, and Kaplan 2024, p. 2356, l. 13--15)}
\item
  \textbf{Controls (W):} gender, age or birth cohort, education, the county share of foreigners, the county unemployment rate, federal state (17), left--right self-placement, religion, being very proud to be German, the view that foreigners enrich German culture and the view that Islam fits into German society.\footnote{(Binstok, Gould, and Kaplan 2024, p. 2352, l. 15--23)}

  \begin{itemize}
  \item
    The controls are motivated by earlier work linking antisemitism and anti-Israel views to political ideology and to prejudice against other out-groups.\footnote{(Binstok, Gould, and Kaplan 2024, p. 2352, l. 23--33)}
  \end{itemize}
\end{itemize}

\par\textbf{ADL Survey}\par

\begin{itemize}
\item
  The 2004 survey covers ten European countries and respondents aged 18 to 64; it contains an antisemitism index based on eleven questions and an anti-Israel index based on four questions.\footnote{(Binstok, Gould, and Kaplan 2024, p. 2366, l. 1--23 and l. 28--29)}
\end{itemize}

